\clearpage
\section*{September 8, 2026: Separate unfinished decisions from unfinished locations}
\addcontentsline{toc}{section}{September 8, 2026: Separate unfinished decisions from unfinished locations}
\textit{Agent-drafted entry.}

Seven existing decisions now reach the current residential candidate file: six records are excluded because they describe older buildings, alterations, additions or conversions, and one duplicate is suppressed while its successor remains. The duplicate decision is applied only after checking that the retained successor still matches the recorded units and building area. Residential records marked for review decline from 197 to 190. No new numerical measurement was invented.

Of the 190 remaining residential review records, 28 have older decisions awaiting verification or application, 158 have no matching decision in the reviewed tables, and four concern the Cirrus and One Chicago conflicts. Another 49 residential records await reconciliation with the commercial source. The list contains 2,346 source records altogether, dominated by the individual-location issue described below. An older decision is not sufficient when later records describe a different object. At One Chicago, for example, the older residential replacement instruction refers to 77 condominiums, while the selected commercial record explicitly counts 735 rental apartments and omits the condominiums. That conflict remains open. Cirrus and Cascade likewise need building-specific treatment; the existing source records must not become three towers.

A separate general location problem affects 2,078 mechanically retained homes: their distances use the centers of older, larger parcels rather than accepted individual home locations. These were already labeled unresolved in the location source field. The new list no longer treats a numerical distance as proof that the individual home has been located. This is a task for a general location rule, not a request for 2,078 manual decisions. Twenty-nine otherwise density-eligible commercial records also lack final locations; settled measurements such as Huron's must be preserved while resolving those locations.

Four earlier commercial recommendations are now applied. Sinclair uses the developer's explicit 2017 completion year. Landmark uses its complete 2021 assessment: 300 apartments, 466,117 building square feet and 39,860 land square feet, excluding the added neighbor. Trio keeps its 100-unit high-rise and uses completed permit 100083121 at 670 W Wayman, which specifically describes that tower rather than the two seven-story buildings. Aqua remains recorded but is excluded from density because its rental-only record does not provide coherent whole-property measurements. Cascade/Cirrus and One Chicago remain open building-definition conflicts. Residential--commercial overlap and unfinished building measurements remain separately identifiable. These source-record counts overlap within sites and are not additive counts of physical buildings to exclude. The 56 completed commercial land decisions are not reopened.

\paragraph{Reproduction.} The release audit produces \texttt{current\_construction\_questions.csv} and its standard report from the current task outputs and explicit decision-table inputs. The audit does not supply production observations. The default construction build, the audit target, and the logbook build are run through Make; the logbook retains saved earlier exhibits. This entry follows commit \texttt{006bbc90}. The paper's frozen input and regression results remain unchanged. Final assembly still has unresolved dependencies and has not been certified as a complete rebuild. An unintended invocation of the root Make entry point was stopped during the permit-summary producer; no density regression ran, and the frozen construction-input checksum remained unchanged.
